\newpage
\thispagestyle{empty}
\chapter*{\centering Project Overview}
\addcontentsline{toc}{chapter}{Project Overview}

Practical lab exams in university computer science programs are harder to manage than they look. Forty students. Two invigilators. Thirty screens. No single tool to watch all of them at once. That is the problem this project addresses.

The AI-Driven Lab Exam Monitoring and Management System — SmartExam — is a twelve-module platform designed specifically for invigilated, machine-bound practical examinations. It covers the full lifecycle: scheduling a session, verifying student identities, locking down workstations, streaming live monitoring telemetry, auto-saving submissions, grading code against test cases, and generating analytical PDF reports.

Existing tools solve fragments of this problem. Respondus LockDown Browser restricts the browser. MOSS checks plagiarism. Canvas manages courses. None of them combine these capabilities with hardware-level device binding, a live invigilator dashboard, and a structured post-exam evidence trail in one workflow designed for physical lab exams.

SmartExam is built across three applications: a WPF desktop client for students (lockdown environment, HWID binding, focus telemetry), a React web panel for teachers and administrators (exam configuration, live monitoring grid, report downloads), and an ASP.NET Core API server backed by a Neon PostgreSQL database.

Development follows Agile with eight four-week sprints. The first two sprints establish authentication and access control. Middle sprints build the exam client, monitoring services, and submission pipeline. The final sprint integrates all twelve modules and prepares the prototype for demonstration.

The expected outcome is a working prototype that gives institutions a single, auditable platform for managing practical lab exams — from pre-exam planning through post-exam review.
